\chapter{Hemorrhoidectomy}

\section*{What Is This Surgery?}
Hemorrhoidectomy is a surgical procedure aimed at the excision of symptomatic hemorrhoidal tissue, primarily indicated for patients with advanced internal or external hemorrhoids that have not responded to conservative treatments or less invasive office-based interventions. This surgery targets the enlarged vascular cushions located in the anal canal and perianal region, which can lead to significant symptoms such as bleeding, prolapse, pain, and hygiene issues when pathologically altered. The procedure involves the careful removal of the diseased tissue along with its associated vascular supply, with the goal of restoring normal anatomy and function while preserving anal continence. Hemorrhoidectomy is recognized as a definitive treatment for severe hemorrhoidal disease, providing long-term symptom relief and enhancing the quality of life for appropriately selected patients.

\section*{Alternative Names}
\begin{itemize}
  \item Excisional Hemorrhoidectomy
  \item Conventional Hemorrhoidectomy
  \item Milligan-Morgan Hemorrhoidectomy (open technique)
  \item Ferguson Hemorrhoidectomy (closed technique)
  \item Stapled Hemorrhoidopexy (Procedure for Prolapse and Hemorrhoids, PPH)
  \item Laser Hemorrhoidoplasty
\end{itemize}

\section*{Relevant Surgical Anatomy}
The anal canal, approximately 3-4 cm in length, extends from the anorectal ring to the anal verge and is lined by distinct types of epithelium: columnar epithelium proximally, transitional epithelium at the dentate line, and anoderm distally. The dentate line serves as a critical anatomical landmark, delineating the transition from visceral to somatic innervation and distinguishing internal from external hemorrhoids. Internal hemorrhoids originate above the dentate line, covered by insensitive columnar or transitional epithelium, while external hemorrhoids arise below the dentate line, covered by sensitive anoderm.

Hemorrhoids consist of normal vascular cushions, which are not varicose veins, comprising a rich arteriovenous plexus, smooth muscle, and connective tissue. Typically, there are three primary internal hemorrhoidal complexes located in the left lateral, right anterior, and right posterior positions, with secondary complexes in between. The arterial supply to these cushions is primarily from the superior rectal artery, a continuation of the inferior mesenteric artery, which descends into the anal canal, forming a dense submucosal plexus. Venous drainage above the dentate line occurs via the superior rectal veins into the portal system, while drainage below the dentate line is via the inferior rectal veins into the systemic circulation (internal pudendal vein). The internal anal sphincter, an involuntary smooth muscle, and the external anal sphincter, a voluntary striated muscle, are crucial for maintaining continence and must be preserved during surgery. Lymphatic drainage from above the dentate line follows the superior rectal vessels to the inferior mesenteric nodes, while below the dentate line, it drains to the inguinal lymph nodes.

\section*{Surgical Principle \& Rationale}
The surgical principle of hemorrhoidectomy revolves around the complete excision of diseased hemorrhoidal complexes, which may be engorged, prolapsed, or thrombosed due to factors such as increased intra-abdominal pressure, straining, and weakening of supporting structures. The operative goal is to remove the affected tissue, including the redundant mucosa, submucosa, and underlying vascular pedicles, while meticulously preserving the integrity of the anal sphincter mechanism and ensuring adequate mucosal bridges to prevent anal stenosis. Ligation of the vascular pedicle is essential to control bleeding and prevent recurrence by interrupting the blood supply to the remaining tissue.

In contrast, stapled hemorrhoidopexy operates on a different principle. This technique involves the circumferential resection of a ring of prolapsed rectal mucosa and submucosa approximately 2-4 cm above the dentate line using a specialized circular stapler. This procedure achieves two main objectives: it mechanically lifts and repositions the prolapsed hemorrhoidal cushions back into their normal anatomical position within the anal canal, and it interrupts the arterial blood supply to the hemorrhoidal plexuses by transecting the feeding vessels, leading to their decompression and atrophy. Both approaches aim to alleviate symptoms by addressing the underlying pathology, with the choice of technique often depending on the type and severity of hemorrhoidal disease, patient preference, and surgeon expertise.

\section*{History \& Surgical Evolution}
The treatment of hemorrhoids has a rich historical background, with references to surgical interventions found in ancient Egyptian papyri and the Hippocratic Corpus, which advocated for ligation and excision. The modern era of hemorrhoidectomy began in the 19th century with the development of more systematic surgical approaches, moving away from crude methods that often resulted in significant morbidity and mortality.

In 1937, Milligan and Morgan introduced their seminal open excisional technique in the United Kingdom, which involved excising the hemorrhoidal tissue and allowing the wounds to heal by secondary intention. This method became widely adopted and remains a gold standard due to its effectiveness. Shortly thereafter, in 1959, Ferguson introduced the closed technique in the United States, which involved primary closure of the mucosal defects after excision, aiming for faster healing and reduced pain. The late 20th century saw further innovations, notably the introduction of stapled hemorrhoidopexy by Dr. Antonio Longo in 1998. This technique offered a less painful alternative for prolapsing internal hemorrhoids by repositioning the tissue and reducing blood flow, rather than direct excision. Concurrently, other minimally invasive approaches, such as Doppler-guided hemorrhoidal artery ligation (DG-HAL) and laser hemorrhoidoplasty, have emerged, reflecting a continuous evolution towards less invasive yet effective treatments with reduced postoperative pain.

\section*{Clinical Indications}
Hemorrhoidectomy is indicated for patients with symptomatic hemorrhoidal disease that has failed to respond to conservative management or less invasive office-based procedures. Specific clinical indications include:

\begin{itemize}
  \item Grade III internal hemorrhoids, characterized by prolapse that requires manual reduction after defecation.
  \item Grade IV internal hemorrhoids, which are permanently prolapsed and irreducible.
  \item Mixed internal and external hemorrhoids causing significant, debilitating symptoms such as persistent bleeding, pain, pruritus, or discomfort.
  \item Large, symptomatic external hemorrhoids or extensive perianal skin tags that interfere with hygiene or cause chronic irritation.
  \item Recurrent or persistent hemorrhoidal symptoms despite adequate trials of conservative management and office-based procedures like rubber band ligation, sclerotherapy, or infrared coagulation.
  \item Acute thrombosed external hemorrhoids causing severe, intractable pain, especially if presenting several days after onset when conservative measures are less effective or if the thrombus is extensive.
  \item Chronic, significant hemorrhoidal bleeding leading to iron-deficiency anemia, unresponsive to other treatments.
  \item Strangulated internal hemorrhoids, which present as an emergency with severe pain, edema, and potential for necrosis, requiring urgent surgical intervention.
\end{itemize}

\section*{Clinical Positioning}
Hemorrhoidectomy is generally considered a secondary or tertiary treatment option for hemorrhoidal disease, reserved for cases where less invasive approaches have failed or are deemed inappropriate due to the severity of the condition. For mild to moderate symptoms (Grade I and II internal hemorrhoids), conservative management involving dietary modifications, increased fiber intake, fluid hydration, and topical agents is the primary approach. If conservative measures are insufficient, office-based procedures such as rubber band ligation, sclerotherapy, or infrared coagulation are typically the next step in the treatment algorithm.

Hemorrhoidectomy becomes a primary definitive treatment for patients presenting with advanced disease, specifically Grade III and IV internal hemorrhoids, large external hemorrhoids, or mixed disease causing significant, debilitating symptoms. It is also considered primary for complications such as acute thrombosed external hemorrhoids with severe pain, or strangulated internal hemorrhoids. In essence, it is part of a step-up approach, becoming the primary choice when the disease burden is high, or when less aggressive interventions have proven ineffective or unsuitable for achieving long-term symptom resolution.

\section*{Surgical Technique \& Procedure}
Hemorrhoidectomy is typically performed under general, spinal, or epidural anesthesia. The patient is usually positioned in the prone jack-knife position, though lithotomy may also be used, to provide optimal exposure of the anal canal and perianal region. The perianal area is prepped and draped in a sterile fashion. A gentle anal dilatation may be performed to facilitate exposure, and an anoscope is inserted to identify the hemorrhoidal complexes.

For excisional hemorrhoidectomy (Milligan-Morgan or Ferguson technique), the surgeon identifies the main hemorrhoidal complexes, usually in the left lateral, right anterior, and right posterior positions.

\begin{enumerate}
  \item An elliptical incision is made around the base of each hemorrhoid, extending from the perianal skin to just above the dentate line.
  \item The hemorrhoidal tissue is then carefully dissected away from the underlying internal anal sphincter muscle, preserving a small bridge of anoderm and mucosa between the excised areas to prevent anal stenosis.
  \item The vascular pedicle at the apex of the hemorrhoid is identified, ligated with absorbable suture (e.g., 2-0 or 3-0 chromic gut or vicryl) to control bleeding, and the hemorrhoidal complex is excised.
  \item In the Milligan-Morgan technique, the wounds are left open to heal by secondary intention, which may reduce postoperative pain and the risk of stenosis.
  \item In the Ferguson technique, the mucosal and skin edges are primarily closed with absorbable sutures, aiming for faster healing and potentially less discharge.
\end{enumerate}

For stapled hemorrhoidopexy, a circular anal dilator is inserted, and a purse-string suture is placed circumferentially in the rectal mucosa approximately 2-4 cm above the dentate line. A specialized circular stapler is then introduced, the purse-string suture is tied around its central rod, and the stapler is fired. This action resects a ring of prolapsed mucosa and submucosa, simultaneously stapling the cut edges together, thereby lifting the hemorrhoidal cushions and interrupting their blood supply. After either procedure, the anal canal is inspected for hemostasis, and a local anesthetic block (e.g., bupivacaine with epinephrine) may be administered for prolonged postoperative pain control. No drains are typically placed.

\section*{Preoperative Workup \& Preparation}
Preoperative preparation for hemorrhoidectomy involves several key steps to ensure patient safety and optimize surgical outcomes. A thorough medical history and physical examination are conducted, including a detailed anorectal examination to confirm the diagnosis, assess sphincter tone, and rule out other pathologies such as anal fissures, fistulas, or malignancy. Laboratory tests, such as a complete blood count, electrolyte panel, and coagulation studies, are routinely performed. Depending on the patient's age and comorbidities, an electrocardiogram (EKG) and chest X-ray may be required, along with medical clearance from a primary care physician or cardiologist.

Specific preparatory steps include:

\begin{itemize}
  \item **Bowel Preparation:** Patients are typically instructed to perform a bowel preparation, such as a Fleet enema or oral laxative (e.g., magnesium citrate or polyethylene glycol) the evening before or morning of surgery, to ensure a clean surgical field and reduce the risk of postoperative fecal impaction.
  
  \item **NPO Status:** Patients must adhere to "nil per os" (NPO) guidelines, usually refraining from food and drink for 6-8 hours prior to surgery, to minimize the risk of aspiration during anesthesia.
  
  \item **Medication Review:** All medications are reviewed. Anticoagulants and antiplatelet agents (e.g., aspirin, warfarin, clopidogrel) often need to be discontinued several days before surgery, under medical supervision, to reduce bleeding risk.
  
  \item **Antibiotics:** Prophylactic intravenous antibiotics (e.g., a single dose of cefazolin) are commonly administered shortly before the incision to reduce the risk of surgical site infection, particularly in colorectal procedures.
  
  \item **Informed Consent:** A comprehensive discussion regarding the procedure, potential risks, benefits, alternatives, and expected postoperative course is held with the patient, and informed consent is obtained.
  
  \item **Postoperative Planning:** Patients are advised on anticipated postoperative pain management strategies, the importance of stool softeners, and dietary modifications.
\end{itemize}

\section*{Contraindications \& Precautions}
**Absolute Contraindications:**

\begin{itemize}
  \item Severe, uncorrected coagulopathy or bleeding disorders that cannot be safely managed preoperatively.
  
  \item Active, untreated rectal or anal malignancy, as the primary focus should be on cancer diagnosis and treatment.
  
  \item Active Crohn's disease or severe ulcerative colitis with significant proctitis, due to a significantly increased risk of poor wound healing, fistula formation, and exacerbation of inflammatory bowel disease.
  
  \item Severe immunosuppression with active perianal infection, which could lead to uncontrolled sepsis and impaired healing.
\end{itemize}

**Relative Contraindications and Precautions:**

\begin{itemize}
  \item Pregnancy: Hemorrhoidectomy is generally deferred until the postpartum period unless symptoms are severe, intractable, and unresponsive to all conservative measures, due to potential risks to both mother and fetus.
  
  \item Portal hypertension with large anorectal varices: This condition significantly increases the risk of severe intraoperative and postoperative bleeding and may require specialized management by a hepatologist.
  
  \item Prior extensive anal surgery: May increase the risk of anal stenosis or compromise sphincter function, requiring careful assessment of residual tissue and sphincter integrity.
  
  \item Significant systemic comorbidities: Uncontrolled diabetes, severe cardiovascular disease, or severe chronic obstructive pulmonary disease (COPD) may increase anesthetic and surgical risks, necessitating thorough preoperative optimization and multidisciplinary team involvement.
  
  \item Acute thrombosed external hemorrhoid presenting very early (within 48-72 hours): Often amenable to simple incision and evacuation in an office setting under local anesthesia, avoiding formal hemorrhoidectomy.
  
  \item Inadequate anal examination: A thorough digital rectal exam and anoscopy are essential preoperatively to rule out other pathologies and assess sphincter tone; if incomplete, further evaluation is needed.
\end{itemize}

\section*{Complications \& Risks}
While generally safe, hemorrhoidectomy carries several potential risks and complications, which can be categorized as general surgical risks and procedure-specific risks.

**General Surgical Risks:**

\begin{itemize}
  \item **Anesthesia Complications:** These can include nausea, vomiting, headache, allergic reactions to medications, and more serious, though rare, cardiac or respiratory events.
  
  \item **Bleeding:** Can occur intraoperatively, immediately postoperatively, or as delayed bleeding (up to 1-2 weeks after surgery) due to sloughing of ligatures or wound dehiscence. This may require re-operation or transfusion.
  
  \item **Infection:** Surgical site infection, cellulitis, or abscess formation in the perianal area, though relatively uncommon due to the rich blood supply and prophylactic antibiotics.
  
  \item **Deep Vein Thrombosis (DVT) and Pulmonary Embolism (PE):** Rare but serious complications, minimized by early ambulation and prophylactic measures.
\end{itemize}

**Procedure-Specific Risks:**

\begin{itemize}
  \item **Severe Pain:** Postoperative pain is common and can be significant, often requiring strong analgesics and multimodal pain management strategies.
  
  \item **Urinary Retention:** Difficulty or inability to urinate is a frequent complication (up to 30\%), often due to pain, local anesthetic effects, or reflex spasm of the bladder sphincter, sometimes requiring temporary catheterization.
  
  \item **Anal Stenosis:** Narrowing of the anal canal due to excessive tissue removal or scarring, which can lead to difficulty with bowel movements and may require further surgical intervention (dilatation or anoplasty).
  
  \item **Fecal Incontinence:** Rare but serious complication, ranging from minor soiling to complete loss of bowel control, usually due to injury to the internal or external anal sphincter muscles or nerve damage during dissection.
  
  \item **Anal Fissure or Fistula Formation:** New development of an anal fissure (a tear in the anoderm) or a fistula-in-ano (an abnormal tract between the anal canal and perianal skin) can occur as a complication of healing or trauma.
  
  \item **Recurrence of Hemorrhoids:** While hemorrhoidectomy is definitive, new hemorrhoids can develop over time, or residual tissue may become symptomatic, particularly if not all complexes were addressed.
  
  \item **Mucosal Prolapse (after stapled hemorrhoidopexy):** If the stapler is fired too low or the purse-string suture is inadequate, there can be persistent or recurrent mucosal prolapse.
  
  \item **Rectovaginal Fistula (extremely rare in females):** An abnormal connection between the rectum and vagina, a devastating complication requiring complex repair.
  
  \item **Perianal Skin Tags:** Residual or new skin tags may form, which are usually asymptomatic but can occasionally cause hygiene issues or discomfort.
\end{itemize}

\section*{Postoperative Recovery Timeline}
Immediately after hemorrhoidectomy, patients are transferred to the Post-Anesthesia Care Unit (PACU) for close monitoring of vital signs, pain levels, and any signs of bleeding. Pain management is a primary focus, often involving a combination of oral or intravenous analgesics, local anesthetic blocks, and sometimes patient-controlled analgesia (PCA). Patients are encouraged to ambulate early to prevent complications like deep vein thrombosis and to aid in bowel function. Urinary retention is a common concern, and patients are monitored for their ability to void spontaneously before discharge.

Upon discharge, typically within 24 hours, patients are provided with detailed instructions for home care. This includes a regimen of oral pain medications (opioids, NSAIDs), stool softeners (e.g., docusate sodium, psyllium fiber supplements), and a high-fiber diet with ample fluid intake to prevent constipation and straining during bowel movements. Warm sitz baths (soaking the anal area in warm water) are highly recommended several times a day, especially after bowel movements, to soothe the area, reduce sphincter spasm, and promote healing. The first bowel movement after surgery can be painful, and patients are advised not to delay it. Most patients experience significant pain for the first few days, which gradually subsides over 1-2 weeks. Return to light activities is usually possible within 1-2 weeks, but heavy lifting and strenuous activities should be avoided for 4-6 weeks to allow for complete wound healing and prevent complications. Full recovery and resolution of symptoms can take 4-6 weeks, with complete wound epithelialization potentially taking longer for open wounds.

\section*{Surgical Findings \& Pathology}
During hemorrhoidectomy, the surgeon meticulously documents the grade, size, and location of the resected hemorrhoidal complexes, noting any associated findings such as thrombosis, ulceration, or extensive skin tags. The presence of any unexpected lesions or induration is also recorded. The excised tissue is invariably sent for histopathological examination. The pathology report will confirm the presence of benign vascular cushions, characterized by dilated submucosal vessels, smooth muscle hypertrophy, and connective tissue proliferation. Crucially, the pathologist will rule out any unexpected findings such as dysplasia, granulomatous inflammation (suggestive of Crohn's disease), or malignancy (e.g., squamous cell carcinoma or adenocarcinoma), which is a vital aspect of the diagnostic workup and ensures appropriate subsequent management if such findings are present.

\section*{Expected Postoperative Course}
A normal or successful postoperative course following a hemorrhoidectomy is characterized by a progressive and complete resolution of the preoperative symptoms that necessitated the surgery. This includes the cessation of bleeding, the absence of prolapsing tissue, and the relief of chronic pain, itching, or discomfort in the anal region. Postoperatively, the surgical wounds in the anal canal should heal without significant complications, typically over a period of 4 to 6 weeks. Patients can expect initial pain, which gradually diminishes, and a return to normal bowel function without excessive straining. The ultimate goal is a healed anal canal, free from the debilitating symptoms of hemorrhoidal disease, allowing the patient to resume a normal, comfortable quality of life with preserved anal continence.

\section*{Postoperative Care \& Follow-up}
Following a hemorrhoidectomy, regular follow-up is essential to monitor healing, address any concerns, and provide guidance for long-term prevention of recurrence.

\begin{itemize}
  \item **Postoperative Clinic Visits:** The first follow-up visit typically occurs 2-6 weeks after surgery. During this visit, the surgeon will examine the anal canal to assess wound healing, check for any complications like stenosis or infection, and ensure symptom resolution. Further visits may be scheduled as needed, particularly for open wounds that require longer healing.
  
  \item **Suture Removal:** If non-absorbable sutures were used for skin closure (less common now), they would be removed at an early follow-up visit. Most sutures used are absorbable and do not require removal.
  
  \item **Dietary and Lifestyle Counseling:** Ongoing advice on a high-fiber diet (25-30g/day), adequate fluid intake (8-10 glasses/day), regular exercise, and avoiding prolonged sitting or straining during defecation is crucial to prevent recurrence and maintain soft stools.
  
  \item **Pelvic Floor Exercises:** While not routinely prescribed, for patients experiencing minor continence issues or persistent discomfort, referral to a pelvic floor physical therapist may be considered.
  
  \item **Activity Resumption:** Patients are guided on a gradual return to normal activities, with advice to avoid heavy lifting or strenuous exercise for 4-6 weeks post-surgery to prevent wound dehiscence or bleeding.
  
  \item **Warning Signs:** Patients are educated on warning signs that warrant immediate medical attention, such as severe or worsening pain unresponsive to medication, fever, chills, purulent discharge, uncontrolled bright red bleeding, inability to pass urine, or persistent difficulty with bowel movements.
\end{itemize}

No specific imaging or laboratory tests are typically required unless complications are suspected.

\section*{Epidemiology}
Hemorrhoidal disease is an exceedingly common condition, affecting a significant portion of the global population. While precise incidence rates are difficult to ascertain due to underreporting and self-treatment, it is estimated that symptomatic hemorrhoids affect approximately 4-10\% of the general population, with a lifetime prevalence as high as 50\%. The peak incidence typically occurs between the ages of 45 and 65 years, though hemorrhoids can affect individuals of any age, including children. There is a slight male predominance in some studies, but pregnancy is a major risk factor for women, often leading to the development or exacerbation of hemorrhoids due to increased intra-abdominal pressure and hormonal changes. Common indications for hemorrhoidectomy, such as chronic bleeding and prolapse, are frequently encountered in surgical practice. Mortality directly attributable to elective hemorrhoidectomy is extremely rare, virtually zero in otherwise healthy individuals. Morbidity, however, is common, with postoperative pain and urinary retention being the most frequent, affecting a substantial percentage of patients, though serious complications like fecal incontinence or anal stenosis are rare, occurring in less than 1-2\% of cases.

\section*{Outcomes \& Prognosis}
The outcomes and prognosis following hemorrhoidectomy are generally excellent, particularly for patients with advanced symptomatic hemorrhoidal disease. Conventional excisional hemorrhoidectomy boasts high success rates, with 85-95\% of patients experiencing complete resolution of their primary symptoms such as bleeding, prolapse, and pain. Functional outcomes are typically very good, with most patients reporting significant improvement in their quality of life and restoration of normal anal function. Survival is not impacted by hemorrhoidectomy itself, as it is a procedure for a benign condition.

Recurrence rates for conventional excisional hemorrhoidectomy are relatively low, estimated to be in the range of 5-10\% over a 5-year period. Stapled hemorrhoidopexy, while associated with less postoperative pain and a faster return to work, may have a slightly higher recurrence rate for bleeding and prolapse compared to excisional techniques, though it remains an effective option for selected patients, particularly those with Grade II-III internal hemorrhoids. Prognostic factors for successful outcomes include appropriate patient selection, meticulous surgical technique, and diligent adherence to postoperative care instructions, including dietary and lifestyle modifications to prevent straining. While complications can occur, the vast majority of patients achieve long-term relief and satisfaction from the procedure, making it a highly effective treatment for severe hemorrhoidal disease.

\section*{References}
\begin{enumerate}
  \item MedlinePlus. Hemorrhoidectomy. U.S. National Library of Medicine; 2024. Available from: \url{https://medlineplus.gov/ency/article/002947.htm}
  
  \item Schwartz's Principles of Surgery. 11th ed. Brunicardi FC, Andersen DK, Billiar TR, et al., editors. McGraw-Hill Education; 2019. Chapter 30: Anus.
  
  \item Sabiston Textbook of Surgery. 21st ed. Townsend CM Jr, Beauchamp RD, Evers BM, Mattox KL, editors. Elsevier; 2021. Chapter 51: Anorectal Disease.
  
  \item American Society of Colon and Rectal Surgeons (ASCRS) Clinical Practice Guidelines for the Management of Hemorrhoids. Dis Colon Rectum. 2018;61(3):284-292.
\end{enumerate}

\clearpage